
\subsubsection{2.1 Weight-Space Quality Prediction}

Unterthiner et al. (2020) established the feasibility of predicting neural network accuracy directly from weight tensors without inference. Using spectral norms, Frobenius norms, and layer-wise statistics, they trained regressors achieving $R^2 \approx$ 0.6--0.7 on CNN model populations. However, their analysis was restricted to convolutional architectures within homogeneous model families.

Eilertsen et al. (2020) extended weight-space analysis to classifier characterization, developing techniques for dissecting the structure of neural network weights. Their methodology provided tools for weight distribution analysis but focused on architectural classification rather than quality prediction.

\subsubsection{2.2 Heavy-Tailed Self-Regularization Theory}

Martin and Mahoney (2021) provided theoretical foundation for weight-based quality assessment through Heavy-Tailed Self-Regularization (HT-SR) theory. Using Random Matrix Theory and the Hill estimator, they showed that well-trained networks develop heavy-tailed weight distributions characterized by power-law exponents $\alpha$. Lower $\alpha$ values correlate with better generalization, reflecting implicit regularization during training.

The WeightWatcher tool \citep{weightwatcher2020} implements this analysis, computing layer-wise $\alpha$ values via singular value decomposition. While the tool supports modern architectures including transformers, systematic validation of HT-SR theory on attention mechanisms had not been conducted prior to this work.

\subsubsection{2.3 Model Zoo Construction}

Schürholt et al. (2022) introduced the model zoo paradigm for weight-space research at scale, constructing diverse populations of trained models for population-level analysis. HuggingFace Model Hub provides an implicit model zoo of unprecedented scale and diversity, though Hub models originate from diverse sources with varying training procedures, creating heterogeneity that may challenge cross-model analysis.

